% Options for packages loaded elsewhere
\PassOptionsToPackage{unicode}{hyperref}
\PassOptionsToPackage{hyphens}{url}
\documentclass[
  ignorenonframetext,
]{beamer}
\newif\ifbibliography
\usepackage{pgfpages}
\setbeamertemplate{caption}[numbered]
\setbeamertemplate{caption label separator}{: }
\setbeamercolor{caption name}{fg=normal text.fg}
\beamertemplatenavigationsymbolsempty
% remove section numbering
\setbeamertemplate{part page}{
  \centering
  \begin{beamercolorbox}[sep=16pt,center]{part title}
    \usebeamerfont{part title}\insertpart\par
  \end{beamercolorbox}
}
\setbeamertemplate{section page}{
  \centering
  \begin{beamercolorbox}[sep=12pt,center]{section title}
    \usebeamerfont{section title}\insertsection\par
  \end{beamercolorbox}
}
\setbeamertemplate{subsection page}{
  \centering
  \begin{beamercolorbox}[sep=8pt,center]{subsection title}
    \usebeamerfont{subsection title}\insertsubsection\par
  \end{beamercolorbox}
}
% Prevent slide breaks in the middle of a paragraph
\widowpenalties 1 10000
\raggedbottom
\AtBeginPart{
  \frame{\partpage}
}
\AtBeginSection{
  \ifbibliography
  \else
    \frame{\sectionpage}
  \fi
}
\AtBeginSubsection{
  \frame{\subsectionpage}
}
\usepackage{iftex}
\ifPDFTeX
  \usepackage[T1]{fontenc}
  \usepackage[utf8]{inputenc}
  \usepackage{textcomp} % provide euro and other symbols
\else % if luatex or xetex
  \usepackage{unicode-math} % this also loads fontspec
  \defaultfontfeatures{Scale=MatchLowercase}
  \defaultfontfeatures[\rmfamily]{Ligatures=TeX,Scale=1}
\fi
\usepackage{lmodern}
\ifPDFTeX\else
  % xetex/luatex font selection
\fi
% Use upquote if available, for straight quotes in verbatim environments
\IfFileExists{upquote.sty}{\usepackage{upquote}}{}
\IfFileExists{microtype.sty}{% use microtype if available
  \usepackage[]{microtype}
  \UseMicrotypeSet[protrusion]{basicmath} % disable protrusion for tt fonts
}{}
\makeatletter
\@ifundefined{KOMAClassName}{% if non-KOMA class
  \IfFileExists{parskip.sty}{%
    \usepackage{parskip}
  }{% else
    \setlength{\parindent}{0pt}
    \setlength{\parskip}{6pt plus 2pt minus 1pt}}
}{% if KOMA class
  \KOMAoptions{parskip=half}}
\makeatother
\usepackage{longtable,booktabs,array}
\newcounter{none} % for unnumbered tables
\usepackage{calc} % for calculating minipage widths
\usepackage{caption}
% Make caption package work with longtable
\makeatletter
\def\fnum@table{\tablename~\thetable}
\makeatother
\setlength{\emergencystretch}{3em} % prevent overfull lines
\providecommand{\tightlist}{%
  \setlength{\itemsep}{0pt}\setlength{\parskip}{0pt}}
\usepackage{bookmark}
\IfFileExists{xurl.sty}{\usepackage{xurl}}{} % add URL line breaks if available
\urlstyle{same}
\hypersetup{
  hidelinks,
  pdfcreator={LaTeX via pandoc}}

\author{\texorpdfstring{}{}}
\date{}

\begin{document}

\section{Scope and related work}\label{scope-and-related-work}

\begin{frame}{Positioning}
\protect\phantomsection\label{positioning}
\textbf{Machine learning for big code} surveys cover naturalness,
representation learning, and task taxonomy {[}@allamanis2018survey{]}.
COGANT contributes \textbf{infrastructure}: a documented IR, pipeline,
and export contract rather than a new benchmark or model architecture.

\textbf{Graph neural networks} unify message-passing frameworks for
relational data {[}@wu2020comprehensive; @scarselli2009graph{]}.
Although COGANT's primary output is the Active Inference Institute's
Generalized Notation Notation, its program graph can be materialised as
optional tensor views compatible with these frameworks (PyG, DGL) so
that graph-neural-network model papers can cite a specific tensor
layout.

\textbf{Code property graphs} and related security-oriented graph
constructions show how to merge syntactic and semantic edges for
analysis {[}@yamaguchi2014modeling{]}. COGANT's program graph is cognate
but emphasizes \textbf{ML export}, confidence scoring, and multi-stage
IR refinement.
\end{frame}

\begin{frame}{Related tool categories}
\protect\phantomsection\label{related-tool-categories}
\begin{itemize}
\tightlist
\item
  \textbf{Compiler and LLVM IRs} --- richer semantics, heavier
  toolchain; COGANT favors lightweight extraction for dataset building.
\item
  \textbf{SCA and linters} --- rule enforcement focus; COGANT focuses on
  graph generation for downstream learning.
\item
  \textbf{Neural code models} (code2vec, CodeBERT, etc.) --- often
  consume tokens or AST paths; COGANT supplies \textbf{explicit program
  graphs} that graph-neural-network-centric research can ingest
  directly.
\end{itemize}
\end{frame}

\begin{frame}[fragile]{Program analysis for ML}
\protect\phantomsection\label{program-analysis-for-ml}
Several systems address the intersection of program analysis and machine
learning that COGANT operates in:

\textbf{code2seq and code2vec} {[}@alon2019code2vec{]} represent
programs as sets of AST paths between terminal nodes. These path-based
representations are effective for method naming and code summarization
but discard the full graph topology that graph-neural-network-based
models exploit. COGANT preserves the complete program graph ---
including control-flow and data-flow edges --- and can export it in
tensor form, enabling downstream architectures that reason over richer
relational structure.

\textbf{Gated Graph Neural Networks (GGNN)} {[}@li2016gated{]}
introduced gated recurrent propagation for graph-structured program
representations, demonstrating strong results on variable misuse
detection and other code tasks. COGANT's export contract (PyG
\texttt{Data} objects with typed \texttt{edge\_index} and
\texttt{edge\_attr}) is directly compatible with GGNN-style models; the
kind and role indices provide the discrete node and edge types that
typed message-passing layers require.

\textbf{CodeQL} (Semmle/GitHub) provides a declarative query language
over relational representations of code. While CodeQL excels at security
analysis with hand-written queries, its outputs are query results rather
than tensor-ready graph bundles. COGANT occupies the complementary
niche: it produces the graph data that learned models consume, and can
ingest CodeQL query results as an additional evidence source feeding the
confidence model.

\textbf{CodeBERT} {[}@feng2020codebert{]} and related pre-trained models
operate at the token level, learning representations from natural
language and code jointly. These models are complementary to COGANT's
graph-centric approach: CodeBERT embeddings could serve as optional node
features in COGANT's Generalized Notation Notation (GNN) export (the
export schema already reserves dimensions for text embeddings as
documented in \texttt{../cogant/docs/GNN\_EXPORT.md}).

\begin{block}{Feature matrix: COGANT vs.~related tools}
\protect\phantomsection\label{feature-matrix-cogant-vs.-related-tools}
The following matrix contrasts COGANT's capabilities with the related
tools discussed in this section. Entries marked ``✓'' indicate
first-class support; ``partial'' indicates limited or indirect support;
``---'' indicates the feature is out of scope for that tool.

\textbf{Table 4. Feature comparison of program-to-model toolchains.}

{\def\LTcaptype{none} % do not increment counter
\begin{longtable}[]{@{}lccccc@{}}
\toprule\noalign{}
Feature & COGANT & code2vec & GGNN & CodeQL & CodeBERT \\
\midrule\noalign{}
\endhead
Full program graph (AST + CFG + DFG) & ✓ & --- & input-only & ✓ & --- \\
Typed node/edge taxonomy & ✓ & --- & partial & ✓ & --- \\
Confidence scoring per assertion & ✓ & --- & --- & --- & --- \\
Provenance tracking & ✓ & --- & --- & partial & --- \\
State-space extraction & ✓ & --- & --- & --- & --- \\
Temporal regime classification & ✓ & --- & --- & --- & --- \\
Dynamic enrichment (coverage, traces) & ✓ & --- & --- & partial & --- \\
Generalized Notation Notation output & ✓ & --- & --- & --- & --- \\
Tensor export (PyG, DGL, HDF5) & ✓ & partial & input-only & --- & --- \\
Pluggable translation rules & ✓ & --- & --- & ✓ & --- \\
Human review loop & ✓ & --- & --- & partial & --- \\
Multi-language front-ends & roadmap & ✓ & --- & ✓ & ✓ \\
\bottomrule\noalign{}
\end{longtable}
}

COGANT is distinct from the other toolchains in three ways: first, it
explicitly models uncertainty through confidence tiers tied to evidence
provenance; second, it produces a structured Active Inference notation
as its primary output rather than an opaque tensor; and third, it
composes static and dynamic evidence in a single pipeline rather than
specializing to one.
\end{block}
\end{frame}

\begin{frame}[fragile]{Active inference and program behavior}
\protect\phantomsection\label{active-inference-and-program-behavior}
The state-space IR in COGANT's pipeline (states, actions, transitions,
observations) shares structural parallels with \textbf{active inference}
formulations {[}@friston2010free{]}, where an agent maintains beliefs
about hidden states and selects actions to minimize prediction error. In
the program analysis context, the ``agent'' is the analysis pipeline
itself: it observes code artifacts, maintains beliefs about program
behavior (the state-space model), and refines those beliefs as new
evidence (dynamic traces, coverage data) arrives.

This connection is analogical: the \texttt{ConfidenceModel} in
\texttt{../cogant/py/cogant/translate/confidence.py} aggregates evidence
and penalties in a way that suggests belief revision, but it is not a
Bayesian posterior. Future work could formalize a tighter link by
casting rule application as variational inference, where a fixpoint
would represent an approximate posterior over program semantics.
\end{frame}

\begin{frame}{Boundaries}
\protect\phantomsection\label{boundaries}
COGANT does not subsume formal verification, interactive theorem
proving, or full interprocedural pointer analysis unless implemented as
explicit future stages. The SPEC marks Rust acceleration and additional
parsers as staged; the manuscript should be read together with that
table for up-to-date scope.
\end{frame}

\begin{frame}[fragile]{Forward compatibility}
\protect\phantomsection\label{forward-compatibility}
Promoting COGANT into
\href{../../../projects/}{\texttt{../../../projects/}} integrates
manuscript PDF rendering with the template's validation gates.
Cross-references in this folder use paths \textbf{relative to these
Markdown files} (for example
\href{../cogant/docs/}{\texttt{../cogant/docs/}}) so links stay stable
when the tree moves.
\end{frame}

\end{document}
